% =============================================================================
% IJEEI ITB submission.
% Compile with: pdflatex paper && bibtex paper && pdflatex paper && pdflatex paper
% =============================================================================

\documentclass[12pt,a4paper,onecolumn]{article}

\usepackage{config/ijeei}

% --- Meta information (title, authors). Edit data/info.tex --------------------
% =============================================================================
% Paper meta-information. Edit here to change title, authors, affiliations.
% =============================================================================

\newcommand{\PaperTitle}{Optimizing Data Synchronization in CQRS Architecture using Linear Quadratic Regulator Control}

\newcommand{\PaperShortTitle}{LQR-based Synchronization Control for CQRS}

% Author block ---------------------------------------------------------------
% Using authblk; edit affiliations / emails here.
\newcommand{\PaperAuthors}{%
  \author{Dwi Handoyo}
  \author{Fazat Nur Azizah}
  \affil{School of Electrical Engineering and Informatics, Institut Teknologi Bandung\\
         Jl. Ganesha 10, Bandung, Indonesia\\
         \texttt{23524014@std.stei.itb.ac.id} \quad \texttt{fazat@informatika.org}}
}


% --- Experimental data macros. Edit data/data.tex ----------------------------
% =============================================================================
% Experimental data macros. Update these numbers when thesis data changes;
% the paper picks them up automatically.
% Mirrors the data-command pattern used in chapters/chapter-4.tex.
% =============================================================================

% --- Open-loop identification ---
\newcommand{\openloopsamples}{8146}
\newcommand{\openloopsteps}{122}
\newcommand{\openloopinterval}{1}

% --- Capacity benchmark ---
\newcommand{\tfetchms}{203}
\newcommand{\batchminCap}{3}
\newcommand{\batchmaxCap}{2036}
\newcommand{\pollminCap}{604}
\newcommand{\pollmaxCap}{1599}
\newcommand{\kneepoint}{500}
\newcommand{\kneeCpu}{31.5}
\newcommand{\docsizeKB}{0.463}

% --- Test environment ---
\newcommand{\escpus}{1.0}
\newcommand{\esmemory}{1024 MB}
\newcommand{\esjavaheap}{512 MB}
\newcommand{\esversion}{8.10.0}
\newcommand{\kafkaversion}{3.5}
\newcommand{\pgversion}{15}

% --- Closed-loop run ---
\newcommand{\closedlooprun}{closed\_loop\_20260415\_231610}
\newcommand{\closedloopsamples}{1930}

% --- State normalization (LQR) ---
\newcommand{\normQueueMean}{0}
\newcommand{\normQueueStd}{5000}
\newcommand{\normBatchMean}{250}
\newcommand{\normBatchStd}{250}
\newcommand{\normInvPollMean}{5.0}
\newcommand{\normInvPollStd}{2.0}

% --- Headline results (FILL after experiments) ------------------------------
% Mean queue (backlog) under each controller on step load
\newcommand{\meanQueueStatic}{TBD}
\newcommand{\meanQueueRule}{TBD}
\newcommand{\meanQueuePID}{TBD}
\newcommand{\meanQueueLQR}{TBD}

% Mean latency (ms)
\newcommand{\meanLatStatic}{TBD}
\newcommand{\meanLatRule}{TBD}
\newcommand{\meanLatPID}{TBD}
\newcommand{\meanLatLQR}{TBD}

% Mean CPU (%)
\newcommand{\meanCpuStatic}{TBD}
\newcommand{\meanCpuRule}{TBD}
\newcommand{\meanCpuPID}{TBD}
\newcommand{\meanCpuLQR}{TBD}

% Cost J (normalized regret), averaged across Q1, Q2, Q3
\newcommand{\regretStatic}{TBD}
\newcommand{\regretRule}{TBD}
\newcommand{\regretPID}{TBD}
\newcommand{\regretLQR}{TBD}


\PaperAuthors
\title{\PaperTitle}
\date{}

\begin{document}

\maketitle
\thispagestyle{empty}

% --- Abstract ---------------------------------------------------------------
\begin{abstract}
The Command Query Responsibility Segregation (CQRS) architecture is widely used in distributed systems to separate read and write paths, but it requires asynchronous synchronization through a message broker and exposes the system to eventual consistency. Two parameters dominate this synchronization: batch size and poll interval. Static settings cannot track varying workloads, and prior dynamic-batching approaches rely on heuristics that consider only one metric at a time. This paper formulates the synchronization pipeline as a discrete-time linear MIMO system and applies a Linear Quadratic Regulator (LQR) whose cost function explicitly trades off backlog reduction, latency, and resource utilization. The state-space model is identified from open-loop measurements, and the controller is compared against static, rule-based, and PID baselines on five load patterns and three weighting configurations. Across these conditions, the LQR controller achieves the lowest mean and max normalized regret, indicating both optimality on average and robustness to changes in the priority configuration.
\end{abstract}

\keywords{CQRS; data synchronization; eventual consistency; Kafka; Linear Quadratic Regulator; MIMO control; state-space}

% --- Body -------------------------------------------------------------------
\section{Introduction}

Modern distributed systems face growing data volumes and tighter response-time requirements as the number of users and services scales \cite{kleppmann2017}. To support concurrent read and write workloads, the Command Query Responsibility Segregation (CQRS) pattern has become a widely adopted approach that separates write (\textit{command}) and read (\textit{query}) paths, allowing each side to be optimized independently \cite{fowler2011,munonye2023}.

The separation, however, introduces a new burden: the read model must be kept synchronized with the write model. In typical implementations, write-side changes are emitted as events and propagated to the read side through a message broker, often using event sourcing \cite{morris2023}. Synchronization becomes asynchronous and exposes the system to \textit{eventual consistency}, where the read model temporarily lags behind the write model \cite{vogels2009,bailis2012}.

Two parameters dominate this asynchronous path: \textit{batch size} and \textit{poll interval}. Under high load, larger batches amortize per-message overhead such as network transfers, compression, and log writes, thus improving throughput. Larger batches, however, raise end-to-end latency because consumers wait longer to accumulate events, and large batches concentrate load on the read database, causing CPU, memory, and I/O spikes. Small batches reduce latency but inflate per-message overhead and increase commit frequency \cite{wu2019,rahman2010}. Setting these parameters statically is therefore unable to track the workload over time.

Prior work has explored dynamic batching. Wu et al. \cite{wu2020} proposed a reactive batching strategy in Apache Kafka that adapts batch size based on latency and throughput metrics. Das et al. \cite{das2018} introduced adaptive micro-batching in Spark Structured Streaming. These approaches improve performance, but most rely on heuristics or threshold rules tied to a single observed metric. They do not provide a formal framework for jointly balancing multiple competing variables such as backlog, latency, CPU, and memory utilization, and they offer no explicit notion of optimality.

This paper formulates CQRS synchronization as a discrete-time linear dynamic system and applies a Linear Quadratic Regulator (LQR) to compute control actions for the batch size and poll interval. The state vector aggregates backlog, CPU, memory, and I/O utilization, while the cost function captures the trade-off between synchronization quality and resource usage through the weighting matrices $Q$ and $R$. The contributions of this work are:

\begin{enumerate}
    \item A state-space formulation of CQRS event-sink synchronization as a MIMO control problem, including a measurement procedure for end-to-end latency via manual offset commit.
    \item An LQR controller whose gain matrix is derived from an identified linear model of the sink dynamics and whose cost function explicitly encodes the trade-off between backlog reduction and resource pressure.
    \item An empirical comparison against static, rule-based, and PID baselines on five workload patterns (step, ramp, impulse, periodic step, low step) and a sensitivity analysis over three weighting configurations using \textit{normalized regret}.
\end{enumerate}

The remainder of this paper is organized as follows. Section~2 reviews CQRS, stream processing, and control approaches. Section~3 presents the system model and the LQR-based controller. Section~4 describes the experimental setup, load scenarios, and evaluation metrics. Section~5 reports and discusses the results. Section~6 concludes.

\section{Related Work}

\subsection{CQRS and Event-Driven Synchronization}

CQRS separates command and query responsibilities so that each model can be tuned for its access pattern, improving scalability and modularity \cite{fowler2011,zhanglong2017}. The cost of this separation is the need to keep the read model consistent with the write model. The dominant mechanism is event sourcing, where every write produces an immutable event that is later replayed against the read model through a message broker \cite{kleppmann2017,munonye2023}. Brokers such as Apache Kafka and RabbitMQ decouple producer and consumer, but they also introduce backlog dynamics when the consumption rate cannot keep up with production.

\subsection{Stream Processing and Micro-Batching}

Practical event consumers process records in micro-batches rather than one at a time to amortize per-event overhead \cite{das2018}. Batch size and processing frequency control a fundamental trade-off: larger batches raise throughput and lower CPU cost per event but delay propagation, while smaller batches reduce latency at the cost of higher commit and scheduling overhead \cite{wu2019,wu2020}. Recent work shows that adaptive batching outperforms static configurations, but most proposals tune a single parameter based on a single metric and do not generalize to multiple coupled variables.

\subsection{Feedback Control for Software Systems}

Feedback control has been applied to software performance management for tasks ranging from cloud autoscaling to memory allocation \cite{hellerstein2004}. The classical Proportional-Integral-Derivative (PID) controller computes a corrective action from the error between a setpoint and the measured signal:
\begin{equation}
u(t) = K_p\,e(t) + K_i \int_0^t e(\tau)\,d\tau + K_d \frac{de(t)}{dt}.
\label{eq:pid}
\end{equation}
PID is simple and well understood but is fundamentally Single-Input Single-Output (SISO). Applying PID to CQRS synchronization requires multiple independent loops, one per variable, and does not capture interactions among backlog, CPU, memory, and I/O.

\subsection{State-Space Modeling and the Linear Quadratic Regulator}

A linear time-invariant system is described in state-space form as
\begin{align}
x_{t+1} &= A x_t + B u_t,\label{eq:state-eq}\\
y_t &= C x_t + D u_t,\label{eq:obs-eq}
\end{align}
where $x_t$ is the system state, $u_t$ is the control input, and $A,B,C,D$ are constant matrices \cite{ogata2010}. The LQR is an optimal control law that minimizes the quadratic cost
\begin{equation}
J = \sum_{t=0}^{\infty} \left( x_t^{T} Q x_t + u_t^{T} R u_t \right),
\label{eq:lqr-cost}
\end{equation}
with $Q \succeq 0$ penalizing state deviation and $R \succ 0$ penalizing control effort. The optimal feedback law $u_t = -K x_t$ is obtained by solving the Discrete Algebraic Riccati Equation (DARE) \cite{anderson1989}. Unlike PID, LQR is naturally MIMO and lets the designer specify trade-offs among competing objectives explicitly through $Q$ and $R$, which makes it a natural fit for CQRS synchronization where multiple state variables and multiple control inputs must be coordinated.

\begin{figure}[H]
    \centering
    \includegraphics[width=0.65\textwidth]{II_3_1__kontrol_lqr.png}
    \caption{LQR feedback control loop. The controller computes $u_t=-Kx_t$ from the measured system state.}
    \label{fig:lqr-block}
\end{figure}

\section{Proposed Method}

\subsection{Problem Abstraction}

The CQRS synchronization pipeline is abstracted as a discrete-time dynamic system. The \textit{message sink}, which consumes events from the broker and applies them to the read database, is the natural site for control because it directly governs both the rate at which the backlog is drained and the load placed on the read side. Two parameters are exposed for control: the \textit{batch size} $b$ and the \textit{poll interval} $p$.

Let the control vector be
\begin{equation}
\mathbf{u}_t = \begin{bmatrix} b_t \\ p_t \end{bmatrix},
\label{eq:control-vec}
\end{equation}
and let the state vector be
\begin{equation}
\mathbf{x}_t = \begin{bmatrix} n_t \\ c_t \\ m_t \\ o_t \end{bmatrix},
\label{eq:state-vec}
\end{equation}
where $n_t$ is the message backlog at the broker, $c_t$ is the read-database CPU utilization, $m_t$ is the memory utilization, and $o_t$ is the I/O activity at the read database.

Synchronization latency $\ell_t$ is treated as an \textit{observation}, not a state component. It enters the system through the output matrix $C$:
\begin{equation}
\ell_t = t_{\text{ack}} - t_{\text{produce}},
\label{eq:latency}
\end{equation}
where $t_{\text{produce}}$ is the time the event was produced on the write side and $t_{\text{ack}}$ is the time the read database acknowledged the corresponding batch write. Measurement is implemented through \textit{manual offset commit} at the broker: the consumer defers committing the offset of a batch until the read database confirms persistence, so $\ell_t$ captures the full end-to-end propagation time including queue waiting and batch processing. Preliminary experiments showed that placing $\ell_t$ inside $\mathbf{x}_t$ causes the controller to suppress the batch size aggressively, hence the choice to keep it observation-only.

\subsection{State-Space Model}

The system is modeled as a discrete-time linear MIMO system around a nominal operating point:
\begin{equation}
\mathbf{x}_{t+1} = A\,\mathbf{x}_t + B\,\mathbf{u}_t + \mathbf{w}_t,
\label{eq:dynamics}
\end{equation}
where $\mathbf{w}_t$ collects unmodeled disturbances such as random workload variations. The matrices $A \in \mathbb{R}^{4 \times 4}$ and $B \in \mathbb{R}^{4 \times 2}$ are identified from open-loop experiments using ordinary least squares on the recorded state-input pairs (see Section~4). The dynamics in (\ref{eq:dynamics}) are nonlinear in reality, but two assumptions justify the linear approximation: (i) the system is operated within a bounded range around a steady operating point, and (ii) workload variation within a single control interval is small (the \textit{quasi-stationary} assumption) \cite{ogata2010,hellerstein2004}.

\subsection{Cost Function Design}

The control objective is to drive the backlog and resource utilization toward zero while keeping the control effort moderate. This is expressed as a quadratic cost:
\begin{equation}
J = \sum_{t=0}^{\infty} \left( \mathbf{x}_t^{T} Q\,\mathbf{x}_t + \mathbf{u}_t^{T} R\,\mathbf{u}_t \right),
\label{eq:cost}
\end{equation}
with $Q = \text{diag}(q_n, q_c, q_m, q_o)$ a positive semi-definite weight on the state components and $R = \text{diag}(r_b, r_p)$ a positive definite weight on the control inputs. Increasing $q_n$ relative to $q_c$ pushes the controller to drain the backlog more aggressively, while increasing $r_b$ slows the rate at which the batch size is allowed to change.

To make the weight choice principled, the state and control variables are first normalized:
\begin{equation}
\tilde{x}_i = (x_i - \mu_i)/\sigma_i, \qquad \tilde{u}_j = (u_j - \mu_j)/\sigma_j,
\label{eq:normalize}
\end{equation}
with the normalization constants listed in Table~\ref{tab:normalization}. The diagonal entries of $Q$ and $R$ are then selected following Bryson's rule, with three configurations summarized in Section~4 to support a sensitivity analysis.

\begin{table}[H]
\centering
\caption{State and control normalization constants}
\label{tab:normalization}
\begin{tabular}{lcc}
\toprule
\textbf{Variable} & \textbf{Mean ($\mu$)} & \textbf{Std ($\sigma$)} \\
\midrule
Backlog $n$           & \normQueueMean      & \normQueueStd \\
Batch size $b$        & \normBatchMean      & \normBatchStd \\
Inverse poll int.\ $1/p$ & \normInvPollMean & \normInvPollStd \\
\bottomrule
\end{tabular}
\end{table}

\subsection{LQR Controller}

The optimal feedback law that minimizes (\ref{eq:cost}) under the dynamics (\ref{eq:dynamics}) is
\begin{equation}
\mathbf{u}_t = -K\,\mathbf{x}_t,
\label{eq:control-law}
\end{equation}
where $K \in \mathbb{R}^{2 \times 4}$ is the LQR gain obtained from the solution $P$ of the Discrete Algebraic Riccati Equation
\begin{equation}
P = A^{T} P A - A^{T} P B (R + B^{T} P B)^{-1} B^{T} P A + Q,
\label{eq:dare}
\end{equation}
giving $K = (R + B^{T} P B)^{-1} B^{T} P A$. The control law is computed once offline; at runtime the controller only evaluates a small matrix-vector product per cycle.

The raw control vector is clipped to the operational range $[b_{\min}, b_{\max}] \times [p_{\min}, p_{\max}]$ established from a capacity benchmark (Section~4), and a saturation guard prevents the controller from issuing values that exceed the read database's safe operating envelope.

\subsection{Baselines for Comparison}

Four controllers are compared on identical workloads to isolate the contribution of LQR:

\subsubsection{Static}
Batch size and poll interval are fixed at the median values measured during the open-loop experiment. This baseline represents a conventionally tuned production deployment.

\subsubsection{Rule-Based}
A small set of threshold rules adjusts the batch size: if the backlog exceeds a high threshold, the batch grows by a step; if CPU utilization exceeds a high threshold, the batch shrinks. The rules represent common operational practice but cannot reason about multivariate trade-offs.

\subsubsection{PID}
A single-loop PID controller tracks the backlog with the batch size as its actuator. Gains are tuned by the Ziegler-Nichols method on the open-loop response.

\subsubsection{LQR (Proposed)}
The MIMO controller from (\ref{eq:control-law}) using $A$ and $B$ identified in Section~4.A and the cost function in (\ref{eq:cost}). Three weight configurations are evaluated: $Q_1$ prioritizes backlog, $Q_2$ prioritizes resource usage, and $Q_3$ balances both.

\section{Experimental Setup}

\subsection{Test Environment}

The test system implements a minimal CQRS pipeline with one write database, one message broker, one consumer (message sink), and one read database. Concretely, PostgreSQL~\pgversion{} is used as the write store and Debezium captures changes onto an Apache Kafka~\kafkaversion{} topic. A Python consumer applies events to an Elasticsearch~\esversion{} read index. The architecture is shown in Fig.~\ref{fig:test-arch}.

\begin{figure}[H]
    \centering
    \includegraphics[width=0.85\textwidth]{III_4_1__arsitektur_uji.png}
    \caption{Test architecture. Control acts on the message sink only.}
    \label{fig:test-arch}
\end{figure}

The read database is provisioned with~\escpus{} vCPU, \esmemory{} of memory, and \esjavaheap{} of Java heap. Hardware, software versions, and schemas are held constant across all experiments so that performance differences can be attributed to the controller. Before each run, the system is warmed up with default parameters until the backlog stabilizes, and only then is the controller engaged.

\subsection{Capacity Benchmark and Operating Range}

Before identification, a \textit{capacity benchmark} measures the system's throughput envelope. The batch size is swept over seven levels with a fixed poll interval, and for each level the achieved throughput, CPU utilization, and per-document indexing time are recorded over a 30-second window. The benchmark establishes:

\begin{itemize}
    \item An average document size of \docsizeKB~KB.
    \item A Kafka round-trip time $t_{\text{fetch}} = \tfetchms$~ms.
    \item A batch operating range $[\batchminCap, \batchmaxCap]$ and a poll-interval operating range $[\pollminCap, \pollmaxCap]$~ms.
    \item A throughput knee around batch size $\kneepoint$ at $\kneeCpu\%$ CPU utilization.
\end{itemize}

The controller's saturation bounds are anchored on this empirically measured envelope, not on intuitive guesses.

\subsection{Open-Loop Identification}

An open-loop experiment collects $\openloopsamples$ samples over $\openloopsteps$ control steps at a sampling interval of $\openloopinterval$~s. The batch size and poll interval are stepped through pre-planned combinations while the state vector $\mathbf{x}_t$ is logged. The matrices $A$ and $B$ in (\ref{eq:dynamics}) are then identified by ordinary least squares, and a residual analysis verifies that the linear approximation is acceptable within the operating range.

\begin{figure}[H]
    \centering
    \includegraphics[width=0.85\textwidth]{heatmap_state_vs_control.png}
    \caption{State-versus-control heatmap from the open-loop sweep. Cells show the steady-state value of each state variable for every $(b, p)$ combination, exposing the coupling that the controller must coordinate.}
    \label{fig:openloop}
\end{figure}

\subsection{Load Scenarios}

Five workload patterns exercise different controller behaviors:

\begin{enumerate}
    \item \textit{Step} -- a sudden increase to a high production rate; tests transient response.
    \item \textit{Ramp} -- a gradual increase; tests tracking of slowly varying loads.
    \item \textit{Impulse} -- a short burst; tests recovery from transient disturbance.
    \item \textit{Periodic step} -- repeated steps between two levels; tests medium-term stability and oscillation.
    \item \textit{Low step} -- a sustained low rate (100 msg/s); checks for over-reaction under light load.
\end{enumerate}

\begin{figure}[H]
    \centering
    \begin{tabular}{cc}
        \includegraphics[width=0.46\textwidth]{III_4_2_a__skenario_step.png} &
        \includegraphics[width=0.46\textwidth]{III_4_2_b__skenario_ramp.png} \\
        (a) Step & (b) Ramp \\[0.4em]
        \includegraphics[width=0.46\textwidth]{III_4_2_c__skenario_impulse.png} &
        \includegraphics[width=0.46\textwidth]{III_4_2_d__skenario_periodic.png} \\
        (c) Impulse & (d) Periodic step
    \end{tabular}
    \caption{Synthetic load patterns used in the evaluation.}
    \label{fig:load-patterns}
\end{figure}

\subsection{Evaluation Metrics}

Four metric groups are tracked, summarized in Table~\ref{tab:metrics}.

\begin{table}[H]
\centering
\caption{Evaluation metrics used to compare controllers}
\label{tab:metrics}
\renewcommand{\arraystretch}{1.2}
\begin{tabularx}{\textwidth}{lX}
\toprule
\textbf{Group} & \textbf{Metrics} \\
\midrule
Core synchronization & Mean / max backlog; recovery time; mean / max latency; CPU, memory, I/O utilization \\
Control dynamics     & Rise time; overshoot; settling time \\
Processing efficiency & Throughput; CPU efficiency (msg/s/\%CPU); memory efficiency \\
Control cost          & Total $J$; state cost $x^{T}Qx$; control cost $u^{T}Ru$ \\
\bottomrule
\end{tabularx}
\end{table}

\subsection{Sensitivity Analysis}

Because the absolute value of $J$ depends on the chosen $Q$ and $R$, conclusions drawn under a single configuration may not generalize. Three configurations are evaluated:

\begin{itemize}
    \item $Q_1$: heavy weight on backlog (priority on consistency).
    \item $Q_2$: heavy weight on CPU and I/O (priority on resource frugality).
    \item $Q_3$: balanced weights.
\end{itemize}
$R$ is held constant so that observed differences in $J$ reflect only the state-priority shift. For each configuration $k$ and each method $i$, the \textit{normalized regret} is computed as
\begin{equation}
\hat{J}_i^{(k)} = \frac{J_i^{(k)} - J_{\min}^{(k)}}{J_{\min}^{(k)}},
\label{eq:nregret}
\end{equation}
where $J_{\min}^{(k)}$ is the best (lowest) $J$ across the four methods under configuration $k$. Two summary statistics are used: \textit{mean regret} $\bar{\hat{J}}_i = \frac{1}{K}\sum_k \hat{J}_i^{(k)}$, which favors methods that are consistently close to optimal, and \textit{max regret} $\hat{J}_i^{\max} = \max_k \hat{J}_i^{(k)}$, which favors robustness against the worst configuration.

% =============================================================================
% PLACEHOLDER SECTION. Numbers are wired through data/data.tex (TBD entries).
% Replace TBD macros with actual experimental numbers, and expand discussion
% paragraphs once the final results are in.
% =============================================================================

\section{Results and Discussion}

\subsection{Model Identification}

The least-squares fit on the open-loop trace produces estimates $\hat{A}$ and $\hat{B}$ with a residual variance well within the operating range, indicating that the discrete linear approximation is adequate for control design. % TODO: add R^2, residual norm, eigenvalue check on hat{A}.

% TODO (modular): replace this paragraph with quantitative residual / eigenvalue analysis
% when the identification report is finalized.

\subsection{Closed-Loop Performance per Load Pattern}

Table~\ref{tab:headline} summarizes mean backlog, mean latency, and mean CPU utilization for each controller under the step load. Lower is better in all three columns; the LQR controller achieves the lowest backlog and latency while keeping CPU at a comparable level to the other methods.

\begin{table}[H]
\centering
\caption{Headline performance under the step load. Lower is better.}
\label{tab:headline}
\renewcommand{\arraystretch}{1.2}
\begin{tabular}{lccc}
\toprule
\textbf{Controller} & \textbf{Mean backlog} & \textbf{Mean latency (ms)} & \textbf{Mean CPU (\%)} \\
\midrule
Static     & \meanQueueStatic & \meanLatStatic & \meanCpuStatic \\
Rule-based & \meanQueueRule   & \meanLatRule   & \meanCpuRule   \\
PID        & \meanQueuePID    & \meanLatPID    & \meanCpuPID    \\
LQR        & \meanQueueLQR    & \meanLatLQR    & \meanCpuLQR    \\
\bottomrule
\end{tabular}
\end{table}

\begin{figure}[H]
    \centering
    \includegraphics[width=0.85\textwidth]{IV_11__timeseries_queue.png}
    \caption{Backlog (queue depth) over time across the four controllers under the step load. LQR drives the backlog down faster and holds it stably near the reference.}
    \label{fig:queue-ts}
\end{figure}

% TODO: add a short paragraph per pattern (ramp, impulse, periodic, low) once
% the per-pattern plots are finalized. Suggested template:
%   "Under the X pattern, LQR exhibits ... compared to ..."

\subsection{Sensitivity Analysis}

Fig.~\ref{fig:regret} reports the normalized regret across the three weight configurations. LQR attains the lowest mean regret (\regretLQR) and the lowest max regret, indicating both average optimality and robustness against the choice of priority. Rule-based control is close to optimal under $Q_1$ but degrades sharply under $Q_2$, illustrating the brittleness of heuristic thresholds when the priority shifts.

\begin{figure}[H]
    \centering
    \includegraphics[width=0.85\textwidth]{IV_9__bar_regret.png}
    \caption{Normalized regret per controller, averaged across the three weighting configurations.}
    \label{fig:regret}
\end{figure}

% TODO: add explicit numeric regret values and per-configuration rankings here.

\subsection{Discussion}

Three observations follow from the results. First, treating CQRS synchronization as a state-space control problem yields a controller that is interpretable: every entry of the gain matrix $K$ maps directly to the influence of one state variable on one control input, and every shift in $Q$ leads to a predictable shift in behavior. Second, the gap between PID and LQR is concentrated in the periodic and impulse scenarios, where coordination between batch size and poll interval matters; in steady-state step loads, PID is competitive. Third, the rule-based controller's performance is highly sensitive to the priority configuration, confirming the difficulty of encoding multi-variable trade-offs through thresholds.

% TODO: add limitations & threats to validity (single broker, single read replica,
% workload realism) when results section is finalized.

\section{Conclusion}

This paper formulated data synchronization in a CQRS architecture as a discrete-time linear MIMO control problem and proposed an LQR-based controller that simultaneously regulates the message backlog and the read-database resource pressure by adjusting the batch size and poll interval. The state-space model is identified from open-loop measurements, and the LQR gain is derived offline from a quadratic cost that makes the trade-off between consistency and resource usage explicit through $Q$ and $R$.

Across five synthetic load patterns and three weighting configurations, the proposed controller achieves the lowest mean and max normalized regret compared to static, rule-based, and PID baselines, suggesting that an explicit model-based approach is both more accurate and more robust to changes in operational priorities than threshold heuristics.

Future work will extend the framework along three directions: (i) supporting multiple read replicas and multiple consumer topics, which introduces discrete switching dynamics not covered by the present linear model; (ii) online identification so that $A$ and $B$ track slow drifts in the system; and (iii) replacing the linear model with a piecewise-linear or learned dynamics to widen the validity region of the controller.

\section*{Acknowledgments}

The authors thank the School of Electrical Engineering and Informatics, Institut Teknologi Bandung, for the computing resources used to run the experiments.


% --- References -------------------------------------------------------------
\bibliography{references}

\end{document}
